%%%%%%%%%%%%%%%%%%%%%%%%%%%%%%%%%%%%%%%%%%%%%%%%%%%%%%%%%%%%%%%%%%%%%%
% Lớp 10 | Hệ thức lượng trong tam giác | Định lý sin và định lý côsin
% Dạng 0H4?2-4: Nhận dạng tam giác
% Số câu: 1  (sinh tự động bởi tach_cau.py; từ nay thầy sửa trực tiếp file này)
%%%%%%%%%%%%%%%%%%%%%%%%%%%%%%%%%%%%%%%%%%%%%%%%%%%%%%%%%%%%%%%%%%%%%%

%%%%%%%%%%%%%%%% TRẮC NGHIỆM 4 PHƯƠNG ÁN %%%%%%%%%%%%%%%%

% nguồn: MaHoaNH_TracNghiem_Lop10
\begin{ex}%[0H4H2-4]%[ma:73382595]
	Tam  giác đều cạnh $a$ nội tiếp trong đường tròn bán kính $R$ bằng
	 \choice
		{$\dfrac{a\sqrt  3}{2}$}
		{\True  $\dfrac{a\sqrt 3}{3}$}
		 {$\dfrac{a\sqrt 2}{3}$}
		{$\dfrac{a\sqrt  3}{4}$}
	\loigiai{ 
		Bán kính đường tròn ngoại tiếp tam giác đều cạnh $a$  thì $R=\dfrac{2}{3}h=\dfrac{2}{3}\cdot\dfrac{a\sqrt 3}{2}=\dfrac{a\sqrt 3}{3}$.
	 }
\end{ex}
